\graphicspath{{abstracts/praesentationen/07-versionierung-archivierung/img/}{img/}}

\begin{beitrag}[Mara Möller, Nadja Neumann]{Versionierung und Archivierung in WissKI-basierten Forschungsprojekten}{%
Mara Möller\,\orcidlink{0000-0001-4567-8901} \\
\small \href{mailto:mara.moeller@example.org}{mara.moeller@example.org} \\
\small Eberhard Karls Universität Tübingen \\[0.5cm]
Nadja Neumann\,\orcidlink{0000-0002-5678-9012} \\
\small \href{mailto:nadja.neumann@example.org}{nadja.neumann@example.org} \\
\small Universitäts- und Landesbibliothek Münster
}

\begin{abstract}
Der Beitrag behandelt Zitierfähigkeit und Langzeitarchivierung von WissKI-Forschungsdaten und unterscheidet drei Versionierungsebenen: redaktionelle Versionierung über Drupal-Revisionen, halbjährliche Schnappschuss-Versionierung als DOI-versehener RDF-Dump in einem institutionellen Repositorium sowie eine abschließende Archivversion als dokumentiertes Submission Information Package gemeinsam mit der Universitätsbibliothek. Empfohlen wird, Versionierung und Archivierung von Beginn an als kontinuierliche Aufgaben einzuplanen und das Archivformat in etablierten RDF-Standards (Turtle, N-Quads) statt in proprietären Triplestore-Exporten zu wählen.
\end{abstract}

\section{Ausgangsfrage}
Forschungsdaten, die in WissKI-Projekten entstehen, sind in der Regel das Ergebnis eines langen, iterativen Erschließungsprozesses. Sie werden über Jahre hinweg ergänzt, korrigiert, restrukturiert. Diese Veränderungsdynamik kollidiert mit zwei wachsenden Anforderungen: Erstens dem Anspruch der \emph{Zitierfähigkeit} -- ein Forschungsbeitrag muss auf einen konkreten, reproduzierbaren Zustand einer Datenbasis verweisen können; zweitens dem Anspruch der \emph{Langzeitarchivierung} -- Forschungsdaten sollen nach Projektende über die Dekaden hinweg verfügbar bleiben. Wie sich diese beiden Anforderungen in WissKI-Projekten operativ umsetzen lassen, ist Gegenstand des vorliegenden Beitrags.

\section{Drei Ebenen der Versionierung}
Wir unterscheiden in unserer Praxis drei Ebenen der Versionierung, die unterschiedliche Aufgaben erfüllen und unterschiedliche Werkzeuge erfordern.

Die erste Ebene ist die \emph{redaktionelle Versionierung}: Jede inhaltliche Änderung an einem WissKI-Datensatz wird über die in Drupal eingebauten Revisionen festgehalten. Diese Ebene ist verlustfrei, aber detailbetont -- sie ist für die Rückverfolgung redaktioneller Entscheidungen geeignet, eignet sich aber nicht als Referenzpunkt für externe Zitationen.

Die zweite Ebene ist die \emph{Schnappschuss-Versionierung}: In definierten Abständen -- in unseren Projekten halbjährlich -- wird der vollständige Triplestore-Inhalt als RDF-Dump exportiert, mit einer DOI versehen und in einem institutionellen Repositorium hinterlegt. Diese Schnappschüsse sind das Mittel der Wahl, um in Publikationen einen konkreten Datenstand zu zitieren. Die Methodik ist andernorts ausführlich diskutiert.\footcite[94]{moeller2024-snapshots}

Die dritte Ebene ist die \emph{Archivversion}: Am Ende eines Projektzyklus -- in der Regel nach Förderende -- wird ein eingefrorener, dokumentierter Datensatz erstellt, der den fachlichen Endstand repräsentiert. Diese Version ist nicht mehr veränderbar und ist mit einer ausführlichen Datendokumentation, einer technischen Dokumentation und einer dauerhaften Lizenz versehen.

\section{Werkzeuge und Workflows}
Für die redaktionelle Versionierung nutzen wir die WissKI-eigenen Drupal-Funktionalitäten ohne nennenswerte Erweiterungen. Für die Schnappschuss-Versionierung haben wir einen Workflow entwickelt, der den Triplestore-Dump zusammen mit einer maschinenlesbaren Schema-Datei, einem READMEt und einem DataCite-konformen Metadatensatz in ein Zenodo-Community-Archiv überführt. Der Workflow ist als Make-getriebene Pipeline implementiert und dokumentiert -- jeder Schnappschuss-Lauf ist damit reproduzierbar und revisionssicher.\footcite[vgl.][128]{neumann2025-archivierung}

Für die Archivversion arbeiten wir eng mit der lokalen Universitätsbibliothek zusammen, die als institutionelles Langzeit-Archiv fungiert. Die Übergabe folgt einem klar definierten Submission Information Package (SIP), das nicht nur die RDF-Daten enthält, sondern auch eine HTML-basierte Sammlungsdokumentation, Screenshots der wichtigsten Oberflächen und eine technische Beschreibung der WissKI-Umgebung zum Zeitpunkt der Archivierung. Diese Begleitdokumentation hat sich als entscheidend erwiesen: Eine in zwanzig Jahren reaktivierte Datenbasis ist ohne diese Begleitdaten praktisch nicht mehr interpretierbar.

\section{Diskussion}
Eine wiederkehrende Diskussion in diesem Themenfeld betrifft das Verhältnis von \emph{lebenden} und \emph{eingefrorenen} Daten. In welchem Maß darf ein WissKI-Projekt nach der Archivierung weiterwachsen? Wir vertreten die Auffassung, dass beide Modi möglich sein sollten, sofern die jeweiligen Zustände klar adressierbar sind: Der archivierte Stand bleibt unverändert, die fortlaufende Pflege erzeugt einen separaten, neu zitierbaren Strang.

Ein weiterer Punkt betrifft das Format der Archivierung selbst. Wir empfehlen ausdrücklich, den RDF-Dump in einem gut etablierten Standardformat (Turtle, N-Quads) bereitzustellen und auf proprietäre Triplestore-Exporte zu verzichten.

\section{Fazit}
Versionierung und Archivierung müssen als kontinuierliche Aufgaben verstanden werden, nicht als Pflichtprogramm am Projektende. Wer diese Aufgaben von Beginn an eingeplant hat, kann am Ende eines Projekts mit überschaubarem Aufwand einen tragfähigen Archivstand abliefern.

\end{beitrag}

\section*{Kurzbiografie(n)}

\subsection*{Mara Möller}
Mara Möller ist Wissenschaftlerin am Tübinger Zentrum für Digital Humanities. Sie beschäftigt sich seit über zehn Jahren mit Versionierungs- und Archivierungsfragen in semantischen Forschungsdatenkontexten und hat mehrere institutionelle Workflows zur Zitierfähigkeit von WissKI-Datenbeständen entwickelt. Aktuell forscht sie zu Standards der Langzeitarchivierung graphbasierter Forschungsdaten.

\subsection*{Nadja Neumann}
Nadja Neumann ist Leiterin der Abteilung Forschungsdatendienste an der Universitäts- und Landesbibliothek Münster. Sie verantwortet die institutionelle Langzeitarchivierung geisteswissenschaftlicher Forschungsdaten und betreut die Übergabe von Forschungsprojekten in den Daueretat der Bibliothek. Ihr Engagement in der bibliothekarischen Standardisierungsarbeit umfasst mehrere Arbeitsgruppen der DINI und der DFG.
